\subsection{Paired domains give a two-thirds-gap separation}
\label{sec:paired-domain}

The isolated point can be farther away than in
Theorem~\ref{cw:punctured-line} with $c_1=c_2=1$. A second construction
uses evaluation points in pairs $x,-x$. Its candidates are even
polynomials, so one successful padding label supplies two agreements.
This construction uses no matching-moment class.

\begin{lemma}[Paired-domain padding]\label{pd:finite}
Let $p>\max\{32,m^2\}$ be prime, $1\le D<m$, and put
\[
 L=\binom mD,\quad U=p-1,\quad R=(p-1)/2-m,\qquad
 B=\frac{(p-1)(p+6\sqrt p+768)}{2p(p-m^2)},\quad
 M=\frac{LR}{R+(L-1)B}.
\]
If $1\le q\le R$ and $U(1-M/U)^q<1$, a length
$n=2(m+q)$ Reed--Solomon code of dimension $K=2D-1$ has a received
line whose every nonzero parameter has at least $2D+2$ agreements.
Parameter zero has maximum agreement exactly $2D$, so there is no
correlated agreement at threshold $2D+2$.
\end{lemma}
\begin{proof}
Choose a tuple $a_1,\ldots,a_m$ of nonzero field elements with distinct
squares, and use the core $\{\pm a_i:1\le i\le m\}$. For each
$D$-subset $I$ of the index set, put
\[
 H_I(X)=\prod_{i\in I}(X^2-a_i^2).
\]
Fix a reference $I_0$, set $w=H_{I_0}$, and take $P_I=w-H_I$.
The leading terms cancel and all odd coefficients vanish, so
$\deg P_I\le2D-2<K$. Each candidate agrees with $w$ on exactly
$2D$ core positions.

We average collisions over all admissible seed tuples. For two supports,
cancel their common factors and let $u=|I\setminus J|$. If $u=1$,
there is no collision outside the core, because the seed squares are
distinct. If $u\ge2$, expose two seed variables $a,b$ occurring only
on one side, and fix the other variables and a nonzero outside point
$x$. The collision equation has the form
\[
 (x^2-a^2)(x^2-b^2)=C,\qquad C\ne0.
\]
This degree-four affine curve is absolutely irreducible. Indeed, over
the algebraic closure, the equation for $a^2$ is
$x^2-C/(x^2-b^2)$, which has simple poles at $b=\pm x$ and hence is
not a square. Its numerator and denominator are coprime, so Gauss's
lemma applies. By \cite[Corollary~5.6]{cafure-matera}, it has at most
$p+6\sqrt p+768$ points.

There are at most $p^{m-2}$ choices for the other seed variables and
$p-1$ choices for $x$. A union bound gives at least
$p^m(1-m^2/p)$ admissible seed tuples: zero entries exclude at most
$mp^{m-1}$ tuples, and equal squares exclude at most
$2\binom m2p^{m-1}$. Divide the collision count by this lower bound
and by two to count the orbits $\{x,-x\}$. The expected number of
outside-orbit collisions for each support pair is at most $B$.
Consequently some core has total collision energy at most
$LR+L(L-1)B$, and its mean outside image size is at least $M$ by
Cauchy--Schwarz. All difference values $P_I(x)-w(x)=-H_I(x)$ are
nonzero and constant on each outside orbit.

Choose the $q$ orbits with largest images. Put $f=w$ everywhere,
$g=0$ on the core, and give each padding pair a shared independent
uniform nonzero value of $g$. The expected number of uncovered nonzero
parameters is at most $U(1-M/U)^q<1$. A choice therefore covers all of
them. A hit supplies two additional agreements. Finally $w$ is monic
of degree $2D$ and the zero codeword already attains its $2D$ roots,
so the exact maximum agreement at parameter zero is $2D$.
\end{proof}

\begin{theorem}[Every point but one, with two-thirds-gap separation]
\label{pd:punctured-line}
Fix a rational rate $\rho\in(0,1)$, $c_1>0$, and $0<c_2<3/2$.
For every sufficiently large prime $p$, a length-$\Theta(\log p)$
Reed--Solomon code of exact rate $\rho$ has a received line such that
all $p-1$ nonzero parameters are within radius
$\theta=1-\rho-\eta$, where $\eta=3/n$, and parameter zero has distance
exactly $\theta+2\eta/3$. The radius is strictly below Elias, there is
no correlated agreement, and
$c_1n2^{c_2H_2(\rho)/\eta}=o(p)$.
\end{theorem}
\begin{proof}
Choose $\rho<\beta<1$ sufficiently close to $\rho$, set
$\alpha=\rho/\beta$, and choose $C$ satisfying
\[
 \frac{2}{\alpha H_2(\beta)}<C<
 \frac{3}{\max\{1,c_2\}H_2(\rho)}.
\]
Take $n=C\log_2p+O(1)$ through odd multiples of the denominator of
$\rho$, and let $K=\rho n$. Set $m=\alpha n/2+O(1)$ and
$D=\rho n/2+O(1)$ with the parity adjustments described below.
Then
\[
 \log_2\binom mD=
 \left(\frac{\alpha C H_2(\beta)}2+o(1)\right)\log_2p,
\]
so $L\ge p^{1+\epsilon}$ for fixed $\epsilon>0$.
Also $UB/R\to1$, hence $M/U\to1$. With
$q=(1-\alpha)n/2+O(1)$, the complete-coverage condition holds.

If $K$ is odd, take $D=(K+1)/2$ and use Lemma~\ref{pd:finite},
adding coordinate zero if $n$ is odd. Give that coordinate values
$f(0)=w(0)$ and $g(0)=0$; no selected candidate agrees there since
$H_I(0)\ne0$. If $K$ is even, the reduced numerator of $\rho$ is
even and its denominator is odd, so $n$ is odd. Take $D=K/2$,
multiply $w$ and all candidates by $X$, and include zero with
$f(0)=g(0)=0$. On a padding pair use directions $xh$ and $-xh$.
The dimension, agreement threshold and exact far agreement each
increase by one compared with the paired construction. In both cases
$A-K=3$, while the far point has maximum agreement $A-2$.
Finally $\eta\log_2p\to3/C>\max\{1,c_2\}H_2(\rho)$ gives
strict Elias and the numerical comparison.
\end{proof}

This theorem strengthens the separation for the usual $c_1=c_2=1$
prescription. Its constant range is explicit; the cubic-moment
construction remains necessary here to defeat arbitrary fixed $c_2$.
